% Part:sets-functions-relations
% Chapter: sets
% Section: schroder-bernstein

\documentclass[../../../include/open-logic-section]{subfiles}

\begin{document}

\olfileid{sfr}{siz}{sb}
\olsection{Het idee van grootte en Schr\"oder-Bernstein}
\begin{explain}
Je zou verwachten dat het zo zit: als $A$ niet groter is dan $B$ en $B$
niet groter is dan $A$, dan hebben $A$ en $B$ evenveel elementen; ze
zijn gelijkmachtig. Eerlijk gezegd: als dat idee \emph{niet klopte},
zou het moeilijk uit te leggen zijn waarom ons begrip van
gelijkmachtigheid, zoals we dat hebben vastgelegd, iets te maken heeft
met het vergelijken van de ‘groottes’ van verzamelingen!{} Gelukkig
klopt dat idee. Dat laat de stelling van Schr\"oder-Bernstein zien.
\end{explain}
\begin{thm}[Schr\"oder-Bernstein]
	\ollabel{thm:schroder-bernstein}
	Als $\cardle{A}{B}$ en $\cardle{B}{A}$,
	dan $\cardeq{A}{B}$.
\end{thm}
\begin{explain}
Anders gezegd: als er !!a{injection} van $A$ naar~$B$ is en
!!a{injection} van $B$ naar~$A$, dan is er !!a{bijection} van $A$
naar~$B$.
Dit is wel \emph{moeilijk} te bewijzen. Dat blijkt ook hieruit: Cantor
formuleerde het resultaat, maar anderen hebben het bewijs
gegeven.\footnote{Lees voor meer over de geschiedenis bijvoorbeeld
\citet[pp.~165--6]{Potter2004}.}
\oliflabeldef{sfr:cardinals:card-sb:sec}{We kunnen pas
\emph{bewijzen} dat Schr\"oder-Bernstein geldt in
\olref[sfr][cardinals][card-sb]{sec}.}{}%
Tot die tijd kun (en moet) je aannemen dat het waar is.
Gelukkig \emph{klopt} Schr\"oder-Bernstein. Dat ondersteunt onze gedachte
dat de relaties $\cardeq{A}{B}$ en $\cardle{A}{B}$ die we hebben
gedefinieerd iets met ‘grootte’ te maken hebben. Ook is
Schr\"oder-Bernstein erg \emph{handig}.
Een !!a{bijection} tussen twee verzamelingen met evenveel elementen is
soms moeilijk te vinden. Dankzij Schr\"oder-Bernstein hoeven we die
bijectie niet in één keer te vinden. We kunnen de twee richtingen apart
aanpakken: eerst zoeken we !!a{injection} van de eerste naar de tweede
verzameling, en daarna hetzelfde in de andere richting.
\end{explain}
\end{document}